\subsection{Exchange sign of the pairing and reality of the auxiliary fields}
\label{p:app:reality}


This appendix checks the exchange sign~\eqref{p:eq:conv:pairing-symmetry} of the pairing on explicit states, with its consequences for the multi-string pairings, and derives the reality of the vertices, of the products and of the auxiliary fields.
We use throughout the normalization~\eqref{p:eq:conv:sphere} and the BPZ rules~\eqref{p:eq:conv:bpz-modes}.


\subsubsection{Exchange sign of the pairing}

The pairing $\pair{A}{B} = \braket{A | c_0^- | B}$ differs from the plain BPZ product $\braket{A | B}$, which is graded-symmetric in the parity of its arguments, by the odd insertion $c_0^-$.
Its exchange sign comes from moving the insertion from one argument to the other.
Moving an operator $O$ from the ket to the bra gives the sign of passing $O$ through $A$:
\begin{equation}
\braket{A | O B}
  = (-1)^{\abs{O} \abs{A}} \, \braket{\transp{O} A | B}.
\end{equation}
Applied to $O = c_0^-$, with $\transp{(c_0^-)} = - c_0^-$ from~\eqref{p:eq:conv:bpz-modes}, we get:
\begin{equation}
  \label{p:eq:rea:exchange-general}
  \begin{aligned}
    \pair{A}{B}
      &
      = \braket{A | c_0^- B}
      = (-1)^{\abs{A} (1 + \abs{B})} \, \braket{c_0^- B | A}
    \\
      &
      = - (-1)^{\abs{A} (1 + \abs{B})} \, (-1)^{\abs{B}} \, \pair{B}{A}
      = (-1)^{(1 + \abs{A}) (1 + \abs{B})} \, \pair{B}{A}.
  \end{aligned}
\end{equation}
The first sign comes from exchanging $A$ with the state $c_0^- B$, of parity $\abs{B} + 1$, in the plain product.
In the second step, we move $c_0^-$ into the bra, which costs the sign $(-1)^{\abs{B}}$ of passing the insertion through $B$ and the minus sign of its BPZ conjugate.
This is~\eqref{p:eq:conv:pairing-symmetry}.
The in-place rule~\eqref{p:eq:conv:bpz-word} brings no further sign: it is the action of the graded anti-involution on a set of creation operators, and a rule reversing the order without a graded sign would be inconsistent beyond the zero modes.

We check~\eqref{p:eq:rea:exchange-general} on the tachyon state $\ket{A} = c_1 \bar c_1 \ket{0}$ and its partner at ghost number three $\ket{B} = c_0^+ c_1 \bar c_1 \ket{0}$.
Their bras follow from~\eqref{p:eq:conv:bpz-modes} and~\eqref{p:eq:conv:bpz-word}, $\bra{A} = \bra{0} c_{-1} \bar c_{-1}$ and $\bra{B} = - \bra{0} c_0^+ c_{-1} \bar c_{-1}$.
  With the normalization~\eqref{p:eq:conv:sphere}, we get:
\begin{equation}
  \label{p:eq:rea:check-values}
  \begin{aligned}
    \pair{A}{B}
      &
      = \braket{c_{-1} \bar c_{-1} c_0^- c_0^+ c_1 \bar c_1}
      = - \frac{1}{2},
    \\
    \pair{B}{A}
      &
      = - \braket{c_0^+ c_{-1} \bar c_{-1} c_0^- c_1 \bar c_1}
      = \braket{c_{-1} \bar c_{-1} c_0^- c_0^+ c_1 \bar c_1}
      = - \frac{1}{2}.
  \end{aligned}
\end{equation}
In the second value, $c_0^+$ passes the even pair $c_{-1} \bar c_{-1}$ without a sign and then anticommutes with $c_0^-$.
The two values agree, as~\eqref{p:eq:conv:pairing-symmetry} requires for an even $A$ and an odd $B$.
  With $\transp{(c_0^\pm)} = + c_0^\pm$, the bra $\bra{B}$ would change sign and the two values would differ.
On the same pair, we have $A^\star = - A$, from the reversal of the two odd modes $c_1 \bar c_1$, and $B^\star = + B$, from the sign $- 1$ of $\transp{(c_0^+)}$ and the reversal sign $- 1$ of three odd modes.
  Hence $\pair{A^\star}{B^\star} = + \tfrac{1}{2} = - \pair{A}{B}^*$, as the equivariance~\eqref{p:eq:conv:pairing-equivariance} requires.

The $n$-string pairing~\eqref{p:eq:conv:pairing-n} inherits the symmetry strand by strand.
We write $X = x_1 \otimes \cdots \otimes x_n$ and $Y = y_1 \otimes \cdots \otimes y_n$.
Exchanging the two tensors exchanges the roles of the $x_i$ and the $y_j$ in the Koszul sign of the definition, and the arguments of every strand pairing:
\begin{equation}
  \label{p:eq:rea:exchange-n}
  \begin{aligned}
    \pair[n]{Y}{X}
      &
      = (-1)^{\sum_{i < j} \abs{x_i} \abs{y_j}} \, \prod_i \pair{y_i}{x_i}
      \\ &
      = (-1)^{\sum_{i < j} \abs{x_i} \abs{y_j} + \sum_i \abs{x_i} \abs{y_i} + \sum_{i < j} \abs{y_i} \abs{x_j}} \, \pair[n]{X}{Y}
      \\ &
      = (-1)^{\abs{X} \abs{Y}} \, \pair[n]{X}{Y}.
  \end{aligned}
\end{equation}
The three sums cover every pair $(i, j)$ exactly once, so the exponent is $\sum_{i, j} \abs{x_i} \abs{y_j} = \abs{X} \abs{Y}$.
This is~\eqref{p:eq:conv:pairing-n-symmetry}, on every pair of homogeneous tensors and by bilinearity on all of $\Hs^{\otimes n}$.

The same bookkeeping fixes the space that the auxiliary fields pair with.
Permuting the strands of both tensors by the same permutation $\pi$, each with the Koszul sign of the physical parity, multiplies the pairing on its support by $\sign(\pi)$.
Indeed, regrouped by pairs of strands, each inversion of $\pi$ exchanges two pairs of total parity $\abs{x_i} + \abs{y_i} = 1$ and costs one sign.
We take $X = \Sigma_n$, with even strands and symmetric components, and $Y$ with strands of ghost number three, all odd.
Then $\pair[n]{\Sigma_n}{Y} = \sign(\pi) \, \pair[n]{\Sigma_n}{\pi Y}$, and only the part of $Y$ with $\pi Y = \sign(\pi) \, Y$ pairs with $\Sigma_n$.
On all-odd strands, the Koszul action of $\pi$ is the plain permutation of the components times $\sign(\pi)$.
The part of $Y$ that pairs with $\Sigma_n$ is thus the part with plainly symmetric components, that is the graded-symmetric tensors in the shifted parity of \cref{p:app:conventions}.

The odd symplectic form~\eqref{p:eq:conv:odd-symplectic}, $\Omega(A, B) = (-1)^{\abs{A}} \, \pair{A}{B}$, is then graded-antisymmetric on every homogeneous pair by~\eqref{p:eq:rea:exchange-general}: the twist $(-1)^{\abs{A}}$ changes sign under the exchange of arguments of opposite parity.


\subparagraph{Reality of the vertices.}

For $n = 2$, the reality~\eqref{p:eq:conv:vertex-reality} follows from the equivariance~\eqref{p:eq:conv:pairing-equivariance} and the star of the BRST operator~\eqref{p:eq:conv:QB-star}:
\begin{equation}
  \label{p:eq:rea:V2-star}
  \begin{aligned}
    \vertex{2}(A^\star, B^\star)
      &
      = \pair{A^\star}{Q_B B^\star}
      = (-1)^{\abs{B} + 1} \, \pair{A^\star}{(Q_B B)^\star}
    \\
      &
      = (-1)^{\abs{B}} \, \pair{A}{Q_B B}^*
      = (-1)^{\abs{A} \abs{B}} \, \vertex{2}(A, B)^*,
  \end{aligned}
\end{equation}
since on the support $\abs{A} = \abs{B}$, so that $\abs{A} \abs{B} = \abs{B}$ modulo two.
Read term by term on the coefficients, with the reality condition~\cite[eq.~(3.6)]{Zwiebach:1993:ClosedStringField} on each of the $n$ fields, the reality of the vertices gives the factor $(-1)^n$ from the sign of that condition and the Koszul factor from reordering the coefficients.
We obtain the reality of the products by stripping the first argument of~\eqref{p:eq:conv:vertex} with~\eqref{p:eq:conv:pairing-equivariance}, since the pairing is non-degenerate on the level-matched space.
In $\vertex{m + 1}(A, B_1, \ldots, B_m) = \pair{A}{\ell_m(B_1, \ldots, B_m)}$, the output has parity $1 + \sum_i \abs{B_i}$.
On the support of the pairing, we thus have $\abs{A} = \sum_i \abs{B_i}$ modulo two, and the rule for arguments of any parity reads:
\begin{equation}
\label{p:eq:rea:product-star}
\ell_m(B_1^\star, \ldots, B_m^\star)
  = (-1)^{m + \sum_i \abs{B_i} + \sum_{i < j} \abs{B_i} \abs{B_j}} \, \ell_m(B_1, \ldots, B_m)^\star.
\end{equation}
On even arguments the sign is $(-1)^m$, and with one odd argument $(-1)^{m + 1}$.

\subparagraph{Reality of the string field.}

The reality condition of the string field~\cite[eq.~(3.6)]{Zwiebach:1993:ClosedStringField}, $\adj{\Psi} = - \bra{\Psi}$, reads $\Psi^\star = - \Psi$ by~\eqref{p:eq:rea:star-bras}.
On the tachyon component
\begin{equation}
\int \frac{\dd^d p}{(2 \pi)^d} \, \phi(p) \, c_1 \bar c_1 \ket{0, p},
\end{equation}
it is the reality condition of a scalar field, $\phi(- p)^* = \phi(p)$, the minus sign coming from the reversal of the two odd modes $c_1 \bar c_1$.
The condition $\Lambda^\star = - \Lambda$ is the one under which every gauge transformation~\eqref{p:eq:conv:gauge} preserves $\Psi^\star = - \Psi$: the free term by the star of $Q_B$,~\eqref{p:eq:conv:QB-star} with $\abs{\Lambda} = 1$, which gives $(Q_B \Lambda)^\star = Q_B \Lambda^\star$ and fixes $\Lambda^\star = - \Lambda$ up to the kernel of $Q_B$, and the interaction terms by the reality of the vertices~\eqref{p:eq:conv:vertex-reality}, through the product rule~\eqref{p:eq:rea:product-star} with one odd argument.

\subparagraph{Reality of the auxiliary fields.}

Every vertex term of the cubic action is a real multiple of a vertex on a tensor whose $N$ strands all have ghost number two.
On even arguments, the reality of the vertices carries the sign $(-1)^N$.
On the slice $\Psi^\star = - \Psi$, $\Sigma_k^\star = (-1)^k \, \Sigma_k$, the strand signs multiply to $(-1)^N$ as well, and we get:
\begin{equation}
\label{p:eq:rea:cubic-real}
\begin{aligned}
  \vertex{N}(\Sigma_k, \Sigma_\ell, \Psi^n)^*
    &
    = (-1)^N \, \vertex{N}\big( \Sigma_k^\star, \Sigma_\ell^\star, (\Psi^\star)^n \big)
  \\
    &
    = (-1)^{N + k + \ell + n} \, \vertex{N}(\Sigma_k, \Sigma_\ell, \Psi^n)
    = \vertex{N}(\Sigma_k, \Sigma_\ell, \Psi^n),
\end{aligned}
\end{equation}
with $N = k + \ell + n$.
The nonpolynomial action is the case without auxiliary fields, $N = n$.
The cubic action is thus real on the slice~\eqref{p:eq:conv:reality-fields} because every strand of ghost number two is star-odd: the field replacing $\Psi^k$ must carry the sign $(-1)^k$ of $(\Psi^\star)^k$.
The sign in the reality condition of $\Sigma_k$ therefore alternates with $k$, and $\Sigma_k$ is not real in a uniform sense.

The source of $\Sigma_k$ is linear in $\Sigma_k$, with a real coefficient once the lower fields are fixed.
It is therefore real exactly when $\Sigma_k$ carries the sign $(-1)^k$ on every direction the source reaches.
For generic vertices, the sources reach every direction and the sign is fixed on the whole field.
For the vertices of a given conformal field theory, the source map need not be surjective.
The directions of $\Sigma_k$ it misses are absent from the dynamics, and on them the condition~\eqref{p:eq:conv:reality-fields} is a choice.

The conditions~\eqref{p:eq:conv:reality-fields} collect these results.
They follow one rule when read on strands: the string field $\Psi$ and the gauge parameter $\Lambda$ are star-odd, and a tensor is star-odd exactly when it has an odd number of star-odd strands.
The strands of ghost number three that occur, those of the outputs of the products on the slice, are star-even, by the same stripping argument as for $\ell_m$.

\subparagraph{Reality of the solution.}

We eliminate the auxiliary fields by inverting their dressed kinetic operators on the sources.
A map built from the vertices commutes with the star on the real slice, by the reality of the vertices, and so does its inverse on a star-stable quotient.
The solution therefore satisfies~\eqref{p:eq:conv:reality-fields} level by level, as a formal series in $\gs$ and wherever the dressed operators are invertible.
Where one of them is degenerate and the source lies in its range, the class of the solution in the quotient is real.
\Cref{p:sec:hs:quintic} works this out at the first order where an inverse appears, with the kernel of $\sprod{2}{2}$ quotiented.
